% ID: SEL-D0055-OB-E01
% ESTADO: validada. S056.
\textbf{a)} $M$ y $N^t$ son de orden $2\times3$, así que se pueden sumar:
\[
M+N^t=\boxed{\begin{pmatrix}2&6&-2\\0&1&-2\end{pmatrix}}.
\]
$M^t$ y $N$ son ambas $3\times2$, así que $M^tN$ no está definido: las dimensiones interiores son $2$ y $3$. En cambio, el producto $(2\times3)(3\times2)$ sí está definido y da
\[
MN=\boxed{\begin{pmatrix}10&3\\4&-1\end{pmatrix}}.
\]

\textbf{b)} Cada entrada suma los productos de los kilos de una fila de $P$ por los precios de una fila de $Q$:
\[
PQ^t=\boxed{\begin{pmatrix}2910&4184\\1372&1972\end{pmatrix}}.
\]
En la diagonal coinciden los tipos de café. Si se vende toda la producción, se obtienen $2910$ unidades monetarias por el café natural y $1972$ por el descafeinado. Por ejemplo, el primer importe es $550\cdot2{,}20+400\cdot2{,}75+240\cdot2{,}50$.
